\begin{definition}
An event $x$ is \emph{guaranteed from} $b$ if every past-inextendible causal
curve from $x$ meets $\Is(b)$. Write $\Det(b)$ for the set of such events.
Clearly $\Is(b)\subseteq\Det(b)$.
\end{definition}

$\Det(b)$ is a property of the spacetime and the event. To connect it with
what the laws fix, let $\mathcal L$ be the set of solutions of some field
equations on $M$. Say that $x$ is \emph{fixed from $b$ under $\mathcal L$},
and write $x\in F_{\mathcal L}(b)$, if any two solutions in $\mathcal L$ that
agree on $\Is(b)$ agree at $x$. The equations, background and globally
admissible solution class are held fixed in this comparison; they are inputs
to the definition, not data an observer is assumed to have measured in the
past. Here agreement means equality of the specified classical fields on the
whole open region, not equality of a finite collection of detector records.
Say that $\mathcal L$ has the
\emph{domain-of-dependence property} if, for every open $U\subseteq M$, any two
solutions that agree on $U$ agree at every event all of whose
past-inextendible causal curves meet $U$. We assume this uniqueness property
explicitly; it is not inferred merely from an operational no-signalling
condition.

\begin{remark}[What is guaranteed and what is fixed]\label{rem:laws}
If $\mathcal L$ has the domain-of-dependence property and consists of
continuous fields, then
\[\Det(b)\cup J^-(b)\subseteq F_{\mathcal L}(b),\]
since $J^-(b)$ is the closure of $\Is(b)$ in a globally hyperbolic spacetime
\citep{MS2008}. So $\Det(b)$ is fixed under every such law: it is what the
causal structure guarantees, whatever the fields. The inclusion is in general
strict, and nothing about $F_{\mathcal L}(b)$ beyond it follows from $\Det(b)$
being small. For example, on Minkowski spacetime with the unit timelike
field $u=\partial_t$, the local equations $u\cdot\nabla Q=P$,
$u\cdot\nabla P=0$ have the solutions $Q=A(\vec X)+tB(\vec X)$, $P=B(\vec X)$
for arbitrary smooth $A,B$. The values of $(Q,P)$ at one point of an integral curve fix them along
it, and the past integral curve from $x$ is a past-inextendible causal curve,
so the law has the domain-of-dependence property; and every integral curve
of $u$ enters $\Is(b)$, so $F_{\mathcal L}(b)=M$ for every $b$. A law of this
kind uses background structure beyond the metric, but nothing in the
domain-of-dependence property excludes it. Neither notion is Belnap's
settledness, on which a proposition is settled at a point event if it holds
in every possible history through it \citep{Belnap1992}; under
indeterministic laws they come apart from it, and we say nothing about such
laws. We use $\Is(b)$, which every local becoming theory counts as real from
$b$, rather than $J^-(b)$; the characterization below uses that $\Is(b)$ is
open.
\end{remark}

\paragraph{Relation to Geroch.} Geroch's domain of prediction $P(q)$ is the set
of $x$ such that every past-directed timelike curve from $x$ without past
endpoint enters $I^-(q)$, but $I^-(x)\not\subseteq I^-(q)$ \citep[sec.~3]{Geroch1977}.
$\Det(b)$ is the causal-curve analogue of the union of $P(b)$ with the
events whose pasts lie in that of $b$; we expect the two to differ at most
on null boundaries, but have not proved it. Geroch gives an equivalent form by achronal surfaces in $I^-(q)$
\citep[sec.~3]{Geroch1977}; Theorem~\ref{thm:char} is a slicing form for
globally hyperbolic spacetimes.

\begin{theorem}\label{thm:char}
Let $t:M\to(\alpha,\beta)$ be a Cauchy time function. Then $x\in\Det(b)$ if
and only if there is $c\in(\alpha,t(x))$ with $J^-(x)\cap S_c\subseteq\Is(b)$.
\end{theorem}

\begin{proof}
($\Leftarrow$) A past-inextendible causal curve from $x$ crosses the Cauchy
surface $S_c$, at a point of $J^-(x)\cap S_c\subseteq\Is(b)$.

($\Rightarrow$) Suppose that for every $c\in(\alpha,t(x))$ the set
$K_c=(J^-(x)\cap S_c)\setminus\Is(b)$ is nonempty; it is compact, since
$\Is(b)$ is open. Fix $t(x)>c_1>c_2>\dots$ with $c_n\to\alpha$, and let
\[P_n=\{(z_1,\dots,z_n)\in K_{c_1}\times\dots\times K_{c_n}:
x\ge z_1\ge z_2\ge\dots\ge z_n\}.\]
$P_n$ is compact, since $\le$ is closed. It is nonempty: pick $y\in K_{c_n}$
and a causal curve from $x$ to $y$; it crosses each $S_{c_i}$, $i<n$, at some
$z_i\le x$, and $z_i\notin\Is(b)$, for otherwise $y\le z_i\ll b$ would give
$y\ll b$. The maps $P_{n+1}\to P_n$ forgetting the last coordinate are
continuous, so the inverse limit of the nonempty compact spaces $P_n$ is
nonempty, giving $x\ge z_1\ge z_2\ge\cdots$ with $z_n\in K_{c_n}$. Join
consecutive terms by causal curves and concatenate. Each segment avoids
$\Is(b)$ by push-up, as before. Along the concatenation $t\to\alpha$, so it has
no past endpoint in $M$: it is a past-inextendible causal curve from $x$
missing $\Is(b)$.
\end{proof}

\begin{corollary}\label{cor:cauchy}
If $\Is(b)$ contains a Cauchy surface $S$, then $\Det(b)=M$.
\end{corollary}

\begin{proof}
If $x\in\Is(S)\cup S$, then $x\le s\ll b$ for some $s\in S$, so $x\in\Is(b)$.
If $x\in\Isp(S)$, extend a past-inextendible causal curve from $x$ to an
inextendible one; it meets $S$, and not to the future of $x$, since then
$s'\ll x\le s$ for some $s,s'\in S$, contradicting achronality. So the curve
from $x$ meets $S\subseteq\Is(b)$.
\end{proof}

Since $\Det(b)$ is defined by causal curves alone, it depends only on the
conformal class of the metric; conformal diagrams decide it.
